%----------------------------------------------------------------------------------------
%	PACKAGES AND THEMES
%----------------------------------------------------------------------------------------

\documentclass{beamer}

\mode<presentation> {%

\usetheme{UNM}

%\setbeamertemplate{footline} % To remove the footer line in all slides uncomment this line
%\setbeamertemplate{footline}[page number] % To replace the footer line in all slides with a simple slide count uncomment this line

%\setbeamertemplate{navigation symbols}{} % To remove the navigation symbols from the bottom of all slides uncomment this line
}

\usepackage{graphicx} % Allows including images
\usepackage{booktabs} % Allows the use of \toprule, \midrule and \bottomrule in tables
\usepackage{pgf}
\usepackage{tikz}
\usepackage{float}
\usepackage{verbatim}

% Setup for the diagrams
\usetikzlibrary{positioning,arrows,automata}
\newcommand{\yslant}{0.5}
\newcommand{\xslant}{-0.6}
%----------------------------------------------------------------------------------------
%	TITLE PAGE
%----------------------------------------------------------------------------------------

\title[Blockchain]{Understanding and using the Blockchain for Credentials} % The short title appears at the bottom of every slide, the full title is only on the title page

\author{Institute for Design \& Innovation} % Your name
\institute[UNM] % Your institution as it will appear on the bottom of every slide, may be shorthand to save space
{%
University of New Mexico \\ % Your institution for the title page
\medskip
\textit{informatics.unm.edu} % Your email address
}
\date{\today} % Date, can be changed to a custom date

\begin{document}

\begin{frame}
\titlepage% Print the title page as the first slide
\end{frame}
\addtobeamertemplate{frametitle}{}{%
  \begin{tikzpicture}[remember picture,overlay]
    \node[anchor=north east,yshift=-1pt] at (current page.north east)
    {\includegraphics[width = 1in]{UNMLogo.png}};
\end{tikzpicture}}

\begin{frame}
\frametitle{Overview} % Table of contents slide, comment this block out to remove it
\tableofcontents % Throughout your presentation, if you choose to use \section{} and \subsection{} commands, these will automatically be printed on this slide as an overview of your presentation
\end{frame}

%----------------------------------------------------------------------------------------
%	PRESENTATION SLIDES
%----------------------------------------------------------------------------------------

%------------------------------------------------
\section{Introduction} % Sections can be created in order to organize your presentation into discrete blocks, all sections and subsections are automatically printed in the table of contents as an overview of the talk
%------------------------------------------------


\begin{frame}
\frametitle{Overview}
This presentation is an overview of the blockchain.
\end{frame}

%----------------------------
% USER CREATION
%----------------------------
\section{Overview}

\begin{frame}
  \frametitle{Description}
  Blockchain is a distributed ledger, storing transactions or arbitrary data
  that may be linked together. The most well-known use of blockchain is for
  Bitcoin and other cryptocurrencies. These blockchains have been used to build
  distributed applications and smart contracts, as Ethereum has done.
\end{frame}

\begin{frame}
  \frametitle{Distributed Ledger}
  The distributed ledger is a public record of all data for the lifetime of the
  database. Anyone who ``mines'' on the blockchain has access to all of these
  records. Therefore, these records are all public. There are ways to make
  information private on the blockchain, however, that is not the focus of this
  presentation.
\end{frame}

\begin{frame}
  \frametitle{Other Distributed Databases}
  Blockchains are distinct from other distributed databases. Each node is a
  peer, unlike in a standard distributed database where the nodes are selected.
  Anybody can join in as a blockchain node. This requires a different security
  model, however, one exists that provides strong consistency.
\end{frame}

\section{Background}
\begin{frame}
  \frametitle{Hash Function}
  A hash function converts an input of any size to an output of any size to an
  output of fixed size. The input cannot be determined based on the output,
  however, the output
  \begin{tabular}{rl}
    \tt \tiny ``The quick brown fox jumped over the lazy dog.'' &
    \tt \tiny ``033939779fb22e05e0837645ec43d98c9a895c615066b30692af97b2f0bb9316'' \\
    \tt \tiny ``The quick brown fox jumped over the lazy doggo.'' &
    \tt \tiny ``a1a7beb72282b4b93100af7fb0b4d9db9845f09d677900439d00126a8ea9cd1a''
  \end{tabular}
\end{frame}

\begin{frame}
  \frametitle{Satisfying a Degree Requirement}

  A degree requirement is satisfied if its {\color{chocolate1}take} and
  {\color{skyblue1}hours} values are satisfied based on its subrequirements, which can
  be courses or other degree requirements. The {\color{chocolate1}take} value is
  simple, just count the subrequirements and see if they meet or exceed the
  value.

\end{frame}

\begin{frame}
  \frametitle{The hours value}
  In order to satisfy the \emph{\color{skyblue1}hours} value, the goal is to
  reduce the subrequirements to a \emph{single element}. In order to do this, we
  use the ${\color{chocolate2}\merge}$ operation on all the subrequirement
  elements. First, we define the absolute value method on a requirement (or
  course):

  \[ {\color{skyblue1}|{\color{orange2}R}|} = {\color{chocolate3}\sum} \left\{
      {\color{butter3}r}_{\color{skyblue1}hours} {\color{aluminium4}\ \mid\ } {\color{butter3}r}
  {\color{chocolate1}\ \in}\ {\color{orange2}R} \right\} \]

  Say: ``Take the sum of the \emph{\color{skyblue1}hours} value of all
  subrequirements of $R$.''

\end{frame}

\begin{frame}
  \frametitle{The ${\color{chocolate2}\merge}$ operator}
  This is the definition of the ${\color{chocolate2}\merge}$ operator:

  \[ {\color{chameleon1}A} {\color{chocolate2}\merge} {\color{red1}B} =
    {\color{chocolate2}(}
    {\color{butter3}subrequirements} =
    {\color{chameleon1}A}_{\color{butter3}sub} {\color{chocolate2}\ \cup\ }
    {\color{red1}B}_{\color{butter3}sub}, \]
    \[ {\color{skyblue1}hours} = {\color{skyblue2}\min}{\color{plum2}\{}
        {\color{skyblue2}\min}{\color{chameleon3}\{}
          {\color{skyblue1}|}{\color{chameleon1}A_{\color{butter3}sub}}
          {\color{chocolate1}\setminus} {\color{red1}B_{\color{butter3}sub}}
          {\color{skyblue1}|},
        {\color{chameleon1}A}_{\color{skyblue1}hours} {\color{chameleon3}\}}
        {\color{chocolate3}\ +\ }\] \[ {\color{skyblue2}\min}{\color{red3}\{}
          {\color{skyblue1}|}{\color{red1}B_{\color{butter3}sub}}
          {\color{chocolate1}\setminus}
          {\color{chameleon1}A_{\color{butter3}sub}} {\color{skyblue1}|},
        {\color{red1}B}_{\color{skyblue1}hours}{\color{red3}\}}
        {\color{chocolate3}\ +\ }
        {\color{skyblue1}|}{\color{chameleon1}A_{\color{butter3}sub}}
        {\color{orange2}\ \cap\ }
      {\color{red1}B_{\color{butter3}sub}}{\color{skyblue1}|},  \] \[
        {\color{chameleon1}A}_{\color{skyblue1}hours} {\color{chocolate3}\ +\ }
    {\color{red1}B}_{\color{skyblue1}hours}
{\color{plum2}\}}{\color{chocolate2})} \]

For the \emph{\color{skyblue1}hours} value, say: ``Take the
{\color{skyblue2}minimum} of the {\color{chocolate3}sum}
of the {\color{skyblue1}hours} of the {\color{butter3}subrequirements} that are
unique to ${\color{chameleon1}A}$ (not in
  ${\color{red1}B}$), and the \emph{\color{skyblue1}hours} value of
  ${\color{chameleon1}A}$. Add this
  to the {\color{chocolate3}sum} of the {\color{skyblue1}hours} of the
  {\color{butter3}subrequirements} that are unique to ${\color{red1}B}$ (not in
  ${\color{chameleon1}A}$), or ${\color{red1}B}$'s \emph{\color{skyblue1}}
  value, whichever is {\color{skyblue2}less}. Add this result to the {\color{chocolate3}sum} of
  the {\color{skyblue1}hours} in the {\color{orange2}intersection} of ${\color{chameleon1}A}$ and
  ${\color{red1}B}$. Take this result, or the {\color{chocolate3}sum} of
  ${\color{chameleon1}A}$'s \emph{\color{skyblue1}hours} value, and
  ${\color{red1}B}$'s \emph{\color{skyblue1}hours} value,
  whichever is {\color{skyblue2}less}.''

\end{frame}

\begin{frame}
  \begin{figure}[h]
    \centering
    \resizebox{\linewidth}{!}{%
      \begin{tikzpicture}
        \node[state, minimum size=1cm, draw=plum2] (A) {};
        \reqnode
        A{\color{plum2}D}{\color{chocolate1}T}{\color{skyblue1}{\color{skyblue1}H}}{\color{chameleon3}G}
        \node[state, minimum size=1cm, below right of=A, node distance=2cm,
        xshift=2cm, draw=butter3] (B) {};
        \reqnode B{\color{butter3}R}{\color{chocolate1}T}{\color{skyblue1}H}{\color{chameleon3}G}
        \node[state, minimum size=1cm, below right of=B, node distance=2cm,
        xshift=0.5cm, draw=butter3] (C) {};
        \reqnode C{\color{butter3}R}{\color{chocolate1}T}{\color{skyblue1}H}{\color{chameleon3}G}
        \node[state, minimum size=1cm, below left of=B, node distance=2cm,
        draw=butter3] (D) {};
        \reqnode D{\color{butter3}R}{\color{chocolate1}T}{\color{skyblue1}H}{\color{chameleon3}G}
        \node[state, minimum size=1cm, below left of=A, node distance=2cm,
        xshift=-2cm, draw=butter3] (E) {};
        \reqnode E{\color{butter3}R}{\color{chocolate1}T}{\color{skyblue1}H}{\color{chameleon3}G}
        \node[state, minimum size=1cm, below left of=E, node distance=2cm,
        xshift=-0.5cm, draw=butter3] (F) {};
        \reqnode F{\color{butter3}R}{\color{chocolate1}T}{\color{skyblue1}H}{\color{chameleon3}G}
        \node[state, minimum size=1cm, below right of=E, node distance=2cm,
        draw=butter3] (G) {};
        \reqnode G{\color{butter3}R}{\color{chocolate1}T}{\color{skyblue1}H}{\color{chameleon3}G}
        \node[state, minimum size=1cm, below left of=F, node distance=1.5cm,
        draw=red2] (H) {};
        \coursenode H{\color{red2}C}{\color{skyblue1}H}
        \node[state, minimum size=1cm, below right of=F, node distance=1.5cm,
        draw=red2] (I) {};
        \coursenode I{\color{red2}C}{\color{skyblue1}H}
        \node[state, minimum size=1cm, below right of=G, node distance=1.5cm,
        draw=red2] (J) {};
        \coursenode J{\color{red2}C}{\color{skyblue1}H}
        \node[state, minimum size=1cm, below left of=G, node distance=1.5cm,
        draw=red2] (K) {};
        \coursenode K{\color{red2}C}{\color{skyblue1}H}
        \node[state, minimum size=1cm, below left of=C, node distance=1.5cm,
        draw=red2] (L) {};
        \coursenode L{\color{red2}C}{\color{skyblue1}H}
        \node[state, minimum size=1cm, below right of=C, node distance=1.5cm,
        draw=red2] (M) {};
        \coursenode M{\color{red2}C}{\color{skyblue1}H}
        \node[state, minimum size=1cm, below left of=D, node distance=1.5cm,
        draw=red2] (N) {};
        \coursenode N{\color{red2}C}{\color{skyblue1}H}
        \node[state, minimum size=1cm, below right of=D, node distance=1.5cm,
        draw=red2] (O) {};
        \coursenode O{\color{red2}C}{\color{skyblue1}H}
        \node[right of=F, node distance=1.65cm] {\color{chocolate2}\Huge$\merge$};
        \node[left of=C, node distance=1.65cm] {\color{chocolate2}\huge$\merge$};

        \path (A) edge (B)
        edge (E)
        (B) edge (C) edge (D)
        (E) edge (F) edge (G)
        (C) edge (L) edge (M)
        (D) edge (N) edge (O)
        (F) edge (H) edge (I)
        (G) edge (J) edge (K);
      \end{tikzpicture}
    }
  \end{figure}

\end{frame}

\begin{frame}
  \begin{figure}
    \resizebox{\linewidth}{!}{%
      \begin{tikzpicture}
        \node[state, minimum size=1cm, draw=plum2] (A) {};
        \reqnode
        A{\color{plum2}D}{\color{chocolate1}T}{\color{skyblue1}{\color{skyblue1}H}}{\color{chameleon3}G}
        \node[state, minimum size=1cm, below right of=A, node distance=2cm,
        xshift=2cm, draw=butter3] (B) {};
        \reqnode B{\color{butter3}R}{\color{chocolate1}T}{\color{skyblue1}H}{\color{chameleon3}G}
        \node[state, minimum size=1cm, below left of=A, node distance=2cm,
        xshift=-2cm, draw=butter3] (E) {};
        \reqnode E{\color{butter3}R}{\color{chocolate1}T}{\color{skyblue1}H}{\color{chameleon3}G}
        \node[state, minimum size=1cm, below left of=F, node distance=1.5cm,
        draw=red2] (H) {};
        \coursenode H{\color{red2}C}{\color{skyblue1}H}
        \node[state, minimum size=1cm, below right of=F, node distance=1.5cm,
        draw=red2] (I) {};
        \coursenode I{\color{red2}C}{\color{skyblue1}H}
        \node[state, minimum size=1cm, below right of=G, node distance=1.5cm,
        draw=red2] (J) {};
        \coursenode J{\color{red2}C}{\color{skyblue1}H}
        \node[state, minimum size=1cm, below left of=G, node distance=1.5cm,
        draw=red2] (K) {};
        \coursenode K{\color{red2}C}{\color{skyblue1}H}
        \node[state, minimum size=1cm, below left of=C, node distance=1.5cm,
        draw=red2] (L) {};
        \coursenode L{\color{red2}C}{\color{skyblue1}H}
        \node[state, minimum size=1cm, below right of=C, node distance=1.5cm,
        draw=red2] (M) {};
        \coursenode M{\color{red2}C}{\color{skyblue1}H}
        \node[state, minimum size=1cm, below left of=D, node distance=1.5cm,
        draw=red2] (N) {};
        \coursenode N{\color{red2}C}{\color{skyblue1}H}
        \node[state, minimum size=1cm, below right of=D, node distance=1.5cm,
        draw=red2] (O) {};
        \coursenode O{\color{red2}C}{\color{skyblue1}H}

        \node[state, minimum size=1cm, below of=E, node distance=1.5cm, draw=butter3] (P) {};
        \reqnode
        P{\color{butter3}R}{\color{chocolate1}-}{\color{skyblue1}\merge}{\color{chameleon3}-}

        \node[state, minimum size=1cm, below of=B, node distance=1.5cm, draw=butter3] (Q) {};
        \reqnode
        Q{\color{butter3}R}{\color{chocolate1}-}{\color{skyblue1}\merge}{\color{chameleon3}-}

        \path (A) edge (B)
        edge (E)
        (Q) edge (L) edge (M)
        (Q) edge (N) edge (O)
        (P) edge (H) edge (I)
        (P) edge (J) edge (K)
        (E) edge (P)
        (B) edge (Q);
      \end{tikzpicture}
    }
  \end{figure}
\end{frame}

\begin{frame}
  \begin{figure}
    \resizebox{\linewidth}{!}{%
      \begin{tikzpicture}
        \node[state, minimum size=1cm, draw=plum2] (A) {};
        \reqnode
        A{\color{plum2}D}{\color{chocolate1}T}{\color{skyblue1}{\color{skyblue1}H}}{\color{chameleon3}G}
        \node[state, minimum size=1cm, below right of=A, node distance=2cm,
        xshift=2cm, draw=butter3] (B) {};
        \reqnode B{\color{butter3}R}{\color{chocolate1}T}{\color{skyblue1}H}{\color{chameleon3}G}
        \node[state, minimum size=1cm, below left of=A, node distance=2cm,
        xshift=-2cm, draw=butter3] (E) {};
        \reqnode E{\color{butter3}R}{\color{chocolate1}T}{\color{skyblue1}H}{\color{chameleon3}G}
        \node[state, minimum size=1cm, below left of=F, node distance=1.5cm,
        draw=red2] (H) {};
        \coursenode H{\color{red2}C}{\color{skyblue1}H}
        \node[state, minimum size=1cm, below right of=F, node distance=1.5cm,
        draw=red2] (I) {};
        \coursenode I{\color{red2}C}{\color{skyblue1}H}
        \node[state, minimum size=1cm, below right of=G, node distance=1.5cm,
        draw=red2] (J) {};
        \coursenode J{\color{red2}C}{\color{skyblue1}H}
        \node[state, minimum size=1cm, below left of=G, node distance=1.5cm,
        draw=red2] (K) {};
        \coursenode K{\color{red2}C}{\color{skyblue1}H}
        \node[state, minimum size=1cm, below left of=C, node distance=1.5cm,
        draw=red2] (L) {};
        \coursenode L{\color{red2}C}{\color{skyblue1}H}
        \node[state, minimum size=1cm, below right of=C, node distance=1.5cm,
        draw=red2] (M) {};
        \coursenode M{\color{red2}C}{\color{skyblue1}H}
        \node[state, minimum size=1cm, below left of=D, node distance=1.5cm,
        draw=red2] (N) {};
        \coursenode N{\color{red2}C}{\color{skyblue1}H}
        \node[state, minimum size=1cm, below right of=D, node distance=1.5cm,
        draw=red2] (O) {};
        \coursenode O{\color{red2}C}{\color{skyblue1}H}
        \node[below of=A, node distance=1.65cm] {\color{chocolate2}\Huge$\merge$};

        \path (A) edge (B)
        edge (E)
        (B) edge (L) edge (M)
        (B) edge (N) edge (O)
        (E) edge (H) edge (I)
        (E) edge (J) edge (K);
      \end{tikzpicture}
    }
  \end{figure}
\end{frame}

\begin{frame}
  \begin{figure}
    \resizebox{\linewidth}{!}{%
      \begin{tikzpicture}
        \node[state, minimum size=1cm, draw=plum2] (A) {};
        \reqnode
        A{\color{plum2}D}{\color{chocolate1}T}{\color{skyblue1}{\color{skyblue1}H}}{\color{chameleon3}G}
        \node[state, minimum size=1cm, below of=A, node distance=1.5cm, draw=butter3] (B) {};
        \reqnode B{\color{butter3}R}{\color{chocolate1}-}{\color{skyblue1}\merge}{\color{chameleon3}-}
        \node[state, minimum size=1cm, below left of=F, node distance=1.5cm,
        draw=red2] (H) {};
        \coursenode H{\color{red2}C}{\color{skyblue1}H}
        \node[state, minimum size=1cm, below right of=F, node distance=1.5cm,
        draw=red2] (I) {};
        \coursenode I{\color{red2}C}{\color{skyblue1}H}
        \node[state, minimum size=1cm, below right of=G, node distance=1.5cm,
        draw=red2] (J) {};
        \coursenode J{\color{red2}C}{\color{skyblue1}H}
        \node[state, minimum size=1cm, below left of=G, node distance=1.5cm,
        draw=red2] (K) {};
        \coursenode K{\color{red2}C}{\color{skyblue1}H}
        \node[state, minimum size=1cm, below left of=C, node distance=1.5cm,
        draw=red2] (L) {};
        \coursenode L{\color{red2}C}{\color{skyblue1}H}
        \node[state, minimum size=1cm, below right of=C, node distance=1.5cm,
        draw=red2] (M) {};
        \coursenode M{\color{red2}C}{\color{skyblue1}H}
        \node[state, minimum size=1cm, below left of=D, node distance=1.5cm,
        draw=red2] (N) {};
        \coursenode N{\color{red2}C}{\color{skyblue1}H}
        \node[state, minimum size=1cm, below right of=D, node distance=1.5cm,
        draw=red2] (O) {};
        \coursenode O{\color{red2}C}{\color{skyblue1}H}

        \path (A) edge (B)
        (B) edge (L) edge (M)
        (B) edge (N) edge (O)
        (B) edge (H) edge (I)
        (B) edge (J) edge (K);
      \end{tikzpicture}
    }
  \end{figure}
\end{frame}


\end{document} 
